
\minitoc

% ========================================================================================================
% Introduction 

Dans ce court chapitre nous allons nous attarder sur la notion d'espérance. 
Jusqu'ici, ce n'était qu'un outil pour décrire la tendance d'une série 
de réalisations d'une variable aléatoire. Nous allons ici nous permettre de la conditionner 
pour en déduire des informations sur notre variable de départ. 

% ========================================================================================================
% Espérances conditionnelles de variables aléatoires discrètes  

\section{Espérances conditionnelles sur un espace d'états discret}

\subsection{Exemple Introductif}

Une entreprise souhaite effectuer un questionnaire de satisfaction en vue d'améliorer ses services. 
Pour chaque client, le résultat du sondage peut être représenté par une variable aléaoite $X$ 
à valeurs dans $E = \{1, \dots, 100\}$. le résultat du sondage sera la moyenne empirique des 
réponses et il donnera une bonne estimation de $ \mathbb{E}(X)$. 

Or, pour permettre de personnaliser un peu plus ses services, l'entreprise souhaite connaître 
les caractéristiques des clients ayant donné le même score au sondage. Par exemple, elle souhaite 
étudier les réponses moyennes des clients en fonction de leur âge, de leur sexe, etc... Par exemple 
si $Y$ est la variable aléatoire qui représente l'âge d'un client, on souhaite obtenir la probabilité 
qu'un client d'âge $y$ attribue la note de $x$. Ce sera donc la quantité suivante : 
    $$ \myP(X = x \mid Y = y) = \frac{ \myP(X = x, Y = y) }{ \myP(Y = y) } $$
que l'on appelle \emph{probabilité conditionnelle} de $x$ sachant $y$. On peut ainsi définir 
une nouvelle mesure de probabilités appelée \emph{probabilité conditionnelle} telle que, pour $y$ 
fixé, on ait : 
    $$ \myP(X = \bullet \mid Y) = \frac{ \myP(X = \bullet, Y = y) }{ \myP(Y = y) } $$ 
    $$ \sum_{x \in E} \myP(X = x \mid Y = y) = 1 $$  
Ainsi, dans notre exemple, la note moyenne des clients d'âge $y$ sera \emph{l'espérance conditionnelle} de 
$X$ sachant $Y = y$ définie par : 
    $$ \mathbb{E}(X \mid Y = y) = \sum_{x \in E} x \myP(X = x \mid Y = y) $$
On peut donc définir une fonction $h : E \longmapsto \overline{\R}$ telle que : 
    $$ h = \mathbb{E}(X \mid Y) \quad \text{i.e} \quad \forall y \in E, h(y) = \mathbb{E}(X \mid Y = y) $$   
À noter que $X$ et $Y$ ne sont pas nécessairement à valeurs dans le même ensemble. 

\begin{example}
    Considérons l'expérience aléatoire suivante. On lance un dé équilibré à 6 faces. Soit $y$ son résultat. 
    On lance ensuite $y$ fois une pièce de monnaie équilibrée. On cherche à déterminer le nombre moyen de 
    fois où on tombe sur Pile. 

    Pour cela, on pose $Y$ une variable aléatoire à valeurs dans $ \llbracket 1, 6 \rrbracket $ et, conditionnellement 
    à $\{Y = y\}$ on définit $X$ la variable aléatoire de loi $ \mathcal{B}(y,p)$ où $p$ est la probabilité de 
    faire Pile. On a alors : 
        $$ \forall k \in \llbracket 0, y \rrbracket , \quad \myP(X = k \mid Y = y) = \binom{k}{y} p^k (1-p)^{y-k} $$ 
    d'où $ \mathbb{E}(X \mid Y = y) = yp$ donc $ \mathbb{E}(X \mid Y) = pY$.    
\end{example}

\begin{prop}[Composition]
    Soient deux variables aléatoires discrètes $X$ et $Y$ on a l'égalité suivante : 
        $$ \mathbb{E}( \mathbb{E}(X \mid Y) ) = \mathbb{E}(X) $$  
\end{prop}

\begin{proof}
    Soient $X$ et $Y$ deux variables aléatoires respectivement à valeurs dans $E$ et $E'$. On a : 
        \begin{align*}
            \mathbb{E}( \mathbb{E}(X \mid Y) ) &= \sum_{y \in E'} \mathbb{E}(X \mid Y = y) \myP(Y = y) \quad \quad \text{(Théorème de transfert)} \\ 
            &= \sum_{y \in E'} \myP(Y = y) \left( \sum_{x \in E} x \myP(X = x \mid Y = y) \right) \\ 
            &= \sum_{y \in E} \sum_{x \in E} x \myP(X = x, Y = y) \\ 
            &= \sum_{x \in E} \sum_{y \in E} x \myP(X = x, Y = y) \\ 
            &= \sum_{x \in E} x \myP(X = x) \quad \quad \text{(Probabilités Totales)} \\
            &= \mathbb{E}(X) 
        \end{align*}
\end{proof}

\begin{proposition}
    Pour reprendre l'exemple précédent, on se donne maintenant $y \in E'$ l'âge d'une personne et 
    on cherche à savoir quelle réponse $x \in E$ cette personne a donnée au sondage. 
    On souhaite donc déterminer $X$ à partir des informations fournies par $Y$. 
    On cherche donc $h(Y)$, une variable aléatoire fonction de $Y$, qui minimise l'erreur quadratique moyenne :
    \[
        \inf_h \mathbb{E}\big[(X - h(Y))^2\big] 
        = \inf_h \left\{ \mathbb{E}(X^2) - 2\,\mathbb{E}[h(Y)X] + \mathbb{E}[h(Y)^2] \right\}.
    \]

    On a, par le théorème de transfert :
    \begin{align*}
        \mathbb{E}\big[h(Y)\mathbb{E}(X|Y)\big] 
        &= \sum_{y \in E'} h(y) \mathbb{E}(X \mid Y = y)\, \mathbb{P}(Y = y) \\ 
        &= \sum_{x \in E}\sum_{y \in E'} x\,h(y)\,\mathbb{P}(X = x, Y = y) \\ 
        &= \mathbb{E}\big[h(Y)X\big].
    \end{align*}

    On en déduit que 
    \[
        \mathbb{E}\big[h(Y)(X - \mathbb{E}[X|Y])\big] = 0.
    \]
    Cette relation peut être interprétée comme une \emph{orthogonalité} : 
    si l’on considère le produit scalaire défini sur les variables aléatoires par 
    \[
        \langle U, V \rangle = \mathbb{E}[UV],
    \]
    alors $X - \mathbb{E}[X|Y]$ est orthogonal à toute variable aléatoire de la forme $h(Y)$. 
    Cette interprétation sera approfondie dans la partie suivante.

    Ainsi :
    \begin{align*}
        \inf_h \Big\{ \mathbb{E}(X^2) - 2\,\mathbb{E}[h(Y)X] + \mathbb{E}[h(Y)^2] \Big\}
        &= \inf_h \Big\{ \mathbb{E}(X^2) - 2\,\mathbb{E}\big[h(Y)\mathbb{E}(X|Y)\big] + \mathbb{E}[h(Y)^2] \Big\} \\
        &= \inf_h \Big\{ \mathbb{E}(X^2) - \mathbb{E}\big[\mathbb{E}(X|Y)^2\big] 
           + \mathbb{E}\big[(h(Y) - \mathbb{E}(X|Y))^2\big] \Big\} \\
        &= \mathbb{E}(X^2) - \mathbb{E}\big[\mathbb{E}(X|Y)^2\big]
           + \inf_h \mathbb{E}\big[(h(Y) - \mathbb{E}(X|Y))^2\big].
    \end{align*}

    Le terme d'espérance est minimal lorsque $h(Y) = \mathbb{E}(X|Y)$, 
    ce qui prouve que l'espérance conditionnelle est la meilleure approximation quadratique de $X$ sachant $Y$.
\end{proposition}


% ========================================================================================================
% Définition de l'espérance conditionnelle

\section{Définition de l'espérence conditionnelle}

Cette section fait appel à beaucoup de connaissances de théorie de la mesure. 
Des rappels peuvent s'imposer. On se place ici dans un espace de variables aléatoires 
de carré intégrable $L^2( \Omega, \mathcal{A}, \myP )$. C'est un espace de Hilbert 
doté d'un produit scalaire : 
    $$ 
        \langle ., \rangle : 
            \begin{cases}
                (L^2( \Omega, \mathcal{A}, \myP ))^2 \longrightarrow \R \\ 
                (X,Y) \longmapsto \mathbb{E}(XY) 
            \end{cases}
    $$ 
et muni d'une norme quadratique : 
    $$ 
        \| .\|_2 : 
            \begin{cases}
                L^2( \Omega, \mathcal{A}, \myP ) \longrightarrow \R \\ 
                X \longmapsto \sqrt{ \mathbb{E}(X^2)}
            \end{cases}
    $$

\subsection{Généralités}

\begin{definition}[Tribu engendrée]
    Soit $( \Omega, \mathcal{A}, \myP)$ un espace de probabilité et $X : \Omega \longrightarrow \R$ une variable 
    aléatoire. On définit la tribu engendrée par $X$, notée $ \sigma(X)$, est la tribu engendrée par les sous-ensembles 
    de $ \Omega$ de la forme $ \{X \leqslant a\}, \forall a \in \R$. 

    \medskip 

    De même, $ \sigma(X)$ contient tout les évènnement de la forme $\{X \in B\}$ avec $B \in \mathcal{B}_\R$. Plus généralement, 
    on peut définir $ \sigma(X_1, \dots, X_n)$ pour une collection de $n$ variables aléatoire $X_1, \dots, X_n$.  
    Intuitivement, $ \sigma(X)$ contient toute l'information associée à la variable $X$. 
\end{definition}

\begin{definition}[Variable aléatoire mesurable par rapport à $X$]
    Soit $(\Omega, \mathcal{A}, \mathbb{P})$ un espace de probabilité et $X : \Omega \to \mathbb{R}$ une variable aléatoire. 
    Une variable aléatoire $Z : \Omega \to \mathbb{R}$ est dite \emph{mesurable par rapport à $X$} si elle est mesurable 
    pour la tribu $\sigma(X)$, c’est-à-dire si :
    \[
        \forall B \in \mathcal{B}_{\mathbb{R}}, \quad Z^{-1}(B) \in \sigma(X).
    \]
\end{definition}

\begin{lemma}
    Soient $Y,W$ deux variables aléatoires sur $( \Omega, \mathcal{A}, \myP)$ telles que $Y$ est à valeurs 
    dans $ (E, \mathcal{E} )$ et $W$ dans $(R, \mathcal{B}_\R)$. $W$ est mesurable par rappport à $ \sigma(Y)$ 
    si et seulement si il existe $f \in \mathcal{M}(E,\R) $ telle que $W = f(Y)$. 
\end{lemma}


\subsection{Espérance conditionnelle}

Nous allons maintenant définir l'espérance conditionnelle comme une projection orthogonale.

\begin{proposition}
    Soient $X, Y \in L^2(\Omega, \mathcal{A}, \mathbb{P})$. 
    On souhaite approcher $X$ par une variable aléatoire mesurable par rapport à $Y$. 
    D'après le lemme précédent, toute variable aléatoire mesurable par rapport à $Y$ 
    est de la forme $h(Y)$ avec une fonction borélienne $h : \mathbb{R} \to \mathbb{R}$.

    On cherche donc à approcher $X$ à partir des variables du sous-espace :
    \[
        \mathcal{H} = \{\, h(Y) \mid h \in L^2(\mathbb{R}, \mathcal{B}_{\mathbb{R}}, \mathbb{P}_Y) \,\}
        \subset L^2(\Omega, \mathcal{A}, \mathbb{P}).
    \]
    La variable aléatoire de $\mathcal{H}$ la plus proche de $X$ 
    est, par définition, le \emph{projeté orthogonal} de $X$ sur $\mathcal{H}$.
\end{proposition}

\begin{definition}[Espérance conditionnelle comme projection]
    Dans le même contexte que précédemment, 
    l'espérance conditionnelle $\mathbb{E}(X \mid Y)$ est le \emph{projeté orthogonal} 
    de $X$ sur le sous-espace $\mathcal{H}$ défini ci-dessus. 
    Elle est caractérisée par la relation d’orthogonalité :
    \[
        \forall h(Y) \in \mathcal{H}, 
        \quad \mathbb{E}\big[(X - \mathbb{E}(X \mid Y))\,h(Y)\big] = 0.
    \]
\end{definition}

On peut étendre cette définition au cas $L^1$ grâce au théorème suivant.

\begin{theorem}[Existence]
    Soit $X \in L^1(\Omega, \mathcal{A}, \mathbb{P})$ et 
    $\mathcal{F}$ une sous-tribu de $\mathcal{A}$. 
    Alors il existe une unique variable aléatoire $Z$, 
    mesurable par rapport à $\mathcal{F}$, telle que :
    \begin{enumerate}[label=\roman*)]
        \item $\mathbb{E}(|Z|) < \infty$, 
        \item pour tout $W \in L^\infty(\Omega, \mathcal{F}, \mathbb{P})$, on a
        \[
            \mathbb{E}(XW) = \mathbb{E}(ZW).
        \]
    \end{enumerate}
    On appelle $Z$ \emph{l'espérance conditionnelle} de $X$ sachant $\mathcal{F}$ 
    et on la note $\mathbb{E}(X \mid \mathcal{F})$.
\end{theorem}

\subsection{Pour les variables à densité}

\begin{rappel}[Loi jointe et lois marginales]
    Soit $(X,Y)$ un couple de variables aléatoires à densité de loi jointe $f_{X,Y}$. 
    La \emph{loi marginale} de $y$ par rapport à $x$ est obtenue par : 
        $$ f_Y(y) = \int_\R f_{X,Y}(x,y) \; dx $$ 
    de même pour la \emph{loi marginale} de $x$ : 
        $$ f_X(x) = \int_\R f_{X,Y}(x,y) \; dy $$
\end{rappel}

\begin{definition}[Probabilité Conditionnelle]
    De même que pour les variables discrètes, on définit la probabilité conditionnelle de $X$ par sa densité : 
        $$ f_{X \mid Y = y} = \frac{f_{X,Y}(x,y)}{f_Y(y)} $$
    Ainsi pour $y \in \R$ fixé tel que  $f_Y(y) > 0$, $f_{X \mid Y = y}$ définit une densité sur $\R$ 
    et on peut lui associer : 
        $$ \mathbb{E}(X \mid Y = y) = \int_\R x f_{X \mid Y = y} (x) \; dx $$ 
\end{definition}

\begin{definition}[Espérance Conditionnelle (Cas continu)]
    Soient deux variables aléatoires $X$ et $Y$ à densité de loi jointe $f_{X,Y}$. 
    On définit l'espérance conditionnelle de $X$ par rapport à $Y$ comme la fonction : 
        $$ h : y \longmapsto \mathbb{E}(X \mid Y = y) = \int_\R x f_{X \mid Y = y} (x) \; dx $$
\end{definition}

\begin{remark}
    La condition d'orthogonalité du théorème précédent se vérifie de la même façon. Pour $ h : \R \longrightarrow \R$ 
    bornée, on a : 
        \begin{align*}
            \mathbb{E}( \mathbb{E}( X \mid Y) h(Y)) &= \int_\R  f_Y(y) h(y) \mathbb{E}(X \mid Y = y) \; dy \\ 
            &= \int_\R f_Y(y) h(y) \left( \int_\R x f_{X \mid Y = y}(x) \; dx \right) \; dy \\ 
            &= \int_{\R^2} x h(y) f_{X \mid Y = y}(x) f_Y(y) \; dx dy \\ 
            &= \int_{\R^2} x h(y) f_{X, Y}(x, y) \; dx dy \\ 
            &= \mathbb{E}(Xh(Y)) 
        \end{align*}
\end{remark}

% ========================================================================================================
% Propriétés de l'espérance conditionnelle

\section{Propriétés de l'espérance conditionnelle}

Voyons à présent quelques propriétés de l'espérance conditionnelle. 
Soit $( \Omega, \mathcal{A}, \myP)$ un espace de probabilités. 

\begin{prop}[Calcul]
    Soit $ \mathcal{F}$ une sous-tribu de $ \mathcal{A}$ et $X \in L^1( \mathcal{A}, \myP)$ une variable aléatoire. 
    On a : 
    \begin{enumerate}[label=\roman*)]
        \item $\mathbb{E}( \bullet \mid \mathcal{F} ) $ est \emph{linéaire}. 
        \item si $X \geqslant 0$ alors $ \mathbb{E}(X \mid \mathcal{F} ) \geqslant 0$ {presque sûrement}. 
        \item si $X$ est $ \mathcal{F}$-mesurable alors $ \mathbb{E}(X \mid \mathcal{F} ) = X $ {presque sûrement}. 
    \end{enumerate}  
\end{prop}

\begin{prop}[Indépendance]
    Si $X$ est \emph{indépendante} de $ \mathcal{F}$ alors : 
        $$ \mathbb{E}(X \mid \mathcal{F} ) = \mathbb{E}(X) $$   
\end{prop}

\begin{remark}
    Soit $Y,X$ deux variables aléatoires sur $( \Omega, \mathcal{A}, \myP)$. 
    On pose $ \mathcal{F} = \sigma(Y)$. Dire que $X$ \emph{est indépendante de} $ \mathcal{F} $
    revient à dire que $X$ est indépendante de $Y$. On aura alors : 
        $$ \mathbb{E}(X h(Y)) \overset{\perp}{=} \mathbb{E}(X) \mathbb{E}(h(Y)) = \mathbb{E}( \mathbb{E}(X) h(Y) ) $$    
\end{remark}

Les théorèmes de convergence (monotone et dominée) de l'intégration sont toujours valables pour l'espérance 
conditionnelle. 

\begin{theorem}[Convergence Monotone]
    Soit $(X_n)_{n \in \N}$ une suite de variables aléatoires et $X$ dans $L^1( \mathcal{A}, \myP)$. 
    Si $X_n \underset{n \to \infty}{\longrightarrow} X$ p.s en croissant, alors : 
        $$ \boxed{\lim_{n \to \infty} \mathbb{E}(X_n \mid \mathcal{F}) = \mathbb{E}(X \mid \mathcal{F} )} \quad \text{p.s}  $$
\end{theorem}

\begin{theorem}[Convergence Dominée]
    Soit $(X_n)_{n \in \N}$ une suite de variables aléatoires et $X$ dans $L^1( \mathcal{A}, \myP)$. 
        $$ X_n \underset{n \to \infty}{\longrightarrow} X \; \text{p.s} \quad \text{{et}} \quad  \sup_n |X_n| \leqslant w \in L^1 $$ 
        $$ \Longrightarrow \mathbb{E}(X_n \mid \mathcal{F} ) \underset{n \to \infty}{\longrightarrow} \mathbb{E}(X \mid \mathcal{F} ) $$
\end{theorem}

\begin{theorem}[Inégalité de Jensen]
    Soit $g : \R \longrightarrow \R$ une fonction \emph{convexe} et $X \in L^1( \mathcal{A}, \myP)$. Alors : 
        $$ \mathbb{E}(g(X) \mid \mathcal{F} ) \geqslant g( \mathbb{E}(X \mid \mathcal{F} ) ) $$ 
\end{theorem}

\begin{proposition}
    Soient $X \in L^1( \mathcal{A}, \myP)$ et $ \mathcal{F}, \mathcal{G} $ deux sous-tribus de 
    $ \mathcal{A}$. Alors : 
    \begin{enumerate}[label={\roman*)}]
        \item Si $ \mathcal{F} \subset \mathcal{G}$ (i.e $ \mathcal{G}$ contient plus d'informations 
            que $ \mathcal{F}$) alors : 
                $$ \mathbb{E}( \mathbb{E}(X \mid \mathcal{G} ) \mid \mathcal{F} ) = \mathbb{E}(X \mid \mathcal{F} )  $$
            \emph{(Puisque $ \mathcal{F}$ contient moins d'information que $ \mathcal{G}$, projeter $X$ sur 
            $ \mathcal{G}$ puis sur $ \mathcal{F}$ revient directement à projeter $X$ sur $ \mathcal{F}$)}. 
        \item Si $ \mathcal{F}$ est indépendante de $ \mathcal{G}$ et $X \perp \mathcal{G}$, alors : 
            $$ \mathbb{E}(X \mid \sigma( \mathcal{F}, \mathcal{G})) = \mathbb{E}(X \mid \mathcal{F}) $$ 
        \emph{(Rajouter une information $ \mathcal{G}$ qui n'a rien à voir avec l'expérience ne permet pas 
        d'améliorer la prédiction $ \mathbb{E}(X \mid \mathcal{F}))$}. 
    \item Soit $Y$ une variable aléatoire mesurable par rapport à $ \mathcal{F}$ et que 
        $ XY \in L^1( \mathcal{A}, \myP)$, alors : 
            $$ \mathbb{E}(XY \mid \mathcal{F}) = Y \mathbb{E}(X \mid \mathcal{F}) $$
        \emph{(Si $Y$ ne dépend que de l'information contenue dans $ \mathcal{F}$, alors 
        sa meilleure estimation $ \mathbb{E}(Y \mid \mathcal{F})$ est elle-même.)}
    \end{enumerate} 
\end{proposition}

% \subsection{Linéarité}

% Pour toutes variables aléatoires $X, Y \in L^1$ et pour tous réels $a, b \in \mathbb{R}$, on a :
% \[
%     \mathbb{E}(aX + bY \mid \mathcal{F}) 
%     = a\,\mathbb{E}(X \mid \mathcal{F}) + b\,\mathbb{E}(Y \mid \mathcal{F}).
% \]
% Cette propriété découle directement de la linéarité de l’espérance classique et du fait que la définition
% de $\mathbb{E}(X \mid \mathcal{F})$ repose sur des égalités d’espérance vis-à-vis de variables mesurables
% par rapport à $\mathcal{F}$.

% \subsection{Monotonicité}

% Si $X, Y \in L^1$ et $X \leq Y$ presque sûrement, alors :
% \[
%     \mathbb{E}(X \mid \mathcal{F}) \leq \mathbb{E}(Y \mid \mathcal{F}) \quad \text{p.s.}
% \]
% Cette propriété est une conséquence du théorème de la convergence monotone et de la positivité de l'espérance.

% \subsection{Caractère de projection}

% Si $X$ est $\mathcal{F}$-mesurable, alors :
% \[
%     \mathbb{E}(X \mid \mathcal{F}) = X.
% \]
% Autrement dit, conditionner une variable déjà déterminée par $\mathcal{F}$ ne modifie pas sa valeur.

% De plus, si $\mathcal{F}_1 \subset \mathcal{F}_2$, on a :
% \[
%     \mathbb{E}(\mathbb{E}(X \mid \mathcal{F}_2) \mid \mathcal{F}_1)
%     = \mathbb{E}(X \mid \mathcal{F}_1),
% \]
% ce qu’on appelle la \emph{propriété de la tour}.

% \subsection{Indépendance}

% Si $X$ est indépendante de $\mathcal{F}$, alors :
% \[
%     \mathbb{E}(X \mid \mathcal{F}) = \mathbb{E}(X).
% \]
% En effet, les informations contenues dans $\mathcal{F}$ n’apportent aucune connaissance supplémentaire sur $X$.

% \subsection{Inégalité de Jensen}

% Si $\varphi : \mathbb{R} \to \mathbb{R}$ est une fonction convexe et $X \in L^1$, alors :
% \[
%     \varphi\big(\mathbb{E}(X \mid \mathcal{F})\big)
%     \leq \mathbb{E}\big(\varphi(X) \mid \mathcal{F}\big)
%     \quad \text{p.s.}
% \]
% Cette propriété généralise l’inégalité de Jensen classique au cadre conditionnel.

% \subsection{Cas particulier : conditionnement par une variable aléatoire}

% Lorsque $\mathcal{F} = \sigma(Y)$, toutes les propriétés précédentes s’écrivent avec la notation abrégée
% $\mathbb{E}(X \mid Y)$, par exemple :
% \[
%     \mathbb{E}(\mathbb{E}(X \mid Y)) = \mathbb{E}(X),
%     \quad \text{et} \quad
%     \mathbb{E}(Xh(Y)) = \mathbb{E}(\mathbb{E}(X \mid Y)h(Y)).
% \]

% \begin{remark}
%     La dernière égalité correspond précisément à la relation d’orthogonalité démontrée plus haut :
%     \[
%         \mathbb{E}\big[(X - \mathbb{E}(X \mid Y))\,h(Y)\big] = 0,
%     \]
%     pour toute fonction bornée $h : \mathbb{R} \to \mathbb{R}$.
% \end{remark}


% ========================================================================================================
% Processus Aléatoires

\section{Processus aléatoires - Filtrations}

Les seuls processus que nous avons vu jusqu'à présent sont des suites de variables aléatoires 
indexées par le temps tels que les chaînes de Markov. 
Nous allons ici essayer de généraliser cette notion en introduisant la notion de \emph{filtration}. 

\begin{definition}[Filtration]\label{def:filtration}
    Soit $( \Omega, \mathcal{A}, \myP)$ un espace probabilisé. Une \emph{filtration} de cet espace 
    est une suite \emph{croissante} de tribus $ ( \mathcal{F}_n)_{n \in \N}$ telles que : 
        $$ \mathcal{F}_1 \subset \mathcal{F}_2 \subset \dots \subset \mathcal{A} $$ 
    Une filtration d'un processus sera donc l'information disponible jusqu'au temps $n$ 
    de ce processus. Sa croissance vient du fait que le processus accumule de l'information 
    au cours de son évolution. 

    \medskip

    On parlera ainsi d'espace de probabilité filtré, noté $ ( \Omega, \mathcal{A}, \mathcal{F}, \myP)$. 
\end{definition}

Dans le cadre précédent de l'études des chaînes de Markov, nous avons étudié des suites $(X_n)_{n \in \N}$ 
à valeurs dans $ (\R, \mathcal{B}_\R)$. On se concentrera donc sur l'étude des filtrations de cet espace mesurable. 

Nous pouvons bien sûr réexprimer les chaînes de Markov à l'aide des filtrations en définissant la notion de 
\emph{filtration naturelle} d'une chaîne de Markov. 

\begin{definition}[Filtration Naturelle]
    Soit $(X_n)_{n \in \N}$ une chaîne de Markov sur $( \Omega, \mathcal{A}, \myP)$ valeurs dans $(E, \mathcal{E})$. Sa filtration naturelle est 
    la suite croissante de sous-tribus $ \mathcal{F}_n$  de $ \mathcal{A}$ telles que : 
        \begin{equation}\label{def:filtration-naturelle}
            \forall n \in \N, \quad \mathcal{F}_n = \sigma(X_0, \dots, X_n)
        \end{equation} 
    dans le cas où $E$ est discret, on aura donc que : 
        $$ \forall n \in \N, \quad \mathcal{F}_n = \sigma(\{X_0 = x_0, \dots X_n = x_n \mid \forall x_0, \dots, x_n \in E\}) $$ 
\end{definition}

\begin{remark}
    Soit $(X_n)_{n \in \N}$ une chaîne de Markov sur un espace d'états discret. On a donc définit 
    sa filtration naturelle comme : 
        $$ \forall n \in \N, \quad \mathcal{F}_n = \sigma(\{X_0 = x_0, \dots X_n = x_n \mid \forall x_0, \dots, x_n \in E\}) $$ 
    D'où $ \forall B \in \mathcal{E}$, par propriété de Markov, on a : 
        $$ \myP(X_{n+1} \in B \mid \mathcal{F}_n) = \myP(X_{n+1} \in B \mid X_n) $$ 
    qui est une variable aléatoire à valeurs dans $ (\R, \mathcal{B}_\R)$ qui dépend des valeurs prises par $ X_n$ 
    et dont la valeur est $\myP(X_{n+1} \in B \mid X_n = x_n)$ si $X_n = x_n$. 

    De plus, pour toute fonction mesurable $h : (E, \mathcal{E}) \to (\mathbb{R}, \mathcal{B}_\mathbb{R})$, 
    l'espérance conditionnelle de $h(X_{n+1})$ sachant $X_n$ est une variable aléatoire $\sigma(X_n)$-mesurable définie 
    par :
        $$ \mathbb{E}[h(X_{n+1}) \mid X_n] = g(X_n) $$
    où la fonction $g : (E, \mathcal{E}) \to (\mathbb{R}, \mathcal{B}_\mathbb{R})$ est donnée, dans le cas discret, par :
        $$ g(x_n) = \sum_{y \in E} h(y)\, \mathbb{P}(X_{n+1} = y \mid X_n = x_n) $$
\end{remark}

De même que précédement, on pourra dire qu'un processus est \emph{adapté} à une filtration... 

\begin{definition}[Processus Adapté - Processus Prévisible]
    Soit $( \Omega, \mathcal{A}, \myP)$ un espace de probabilité et $( \mathcal{F}_n)_{n \in \N}$ une filtration de cet espace. 
    Soit $(Y_n)_{n \in \N}$ un processus à valeurs dans $(E, \mathcal{E})$. 
    \begin{enumerate}[label={\roman*)}]
        \item  On dira que le processus $(Y_n)_{n \in \N}$ est \emph{adapté} à la filtration 
        $( \mathcal{F}_n)_{n \in \N}$ si pour tout $n \in \N$, $Y_n$ est $ \mathcal{F}_n$ mesurable. 
                $$ \text{i.e} \quad \forall n \in \N, \quad \forall A \in \mathcal{E}, \quad Y_n^{-1}(A) \in \mathcal{F}_n $$
        \item On dira que le processus $(Y_n)_{n \in \N}$ est \emph{prévisible} par la filtration 
        $( \mathcal{F}_n)_{n \in \N}$ si pour tout $n \in \N$, $Y_n$ est $ \mathcal{F}_{n-1}$ mesurable. 
            $$ \text{i.e} \quad \forall n \in \N, \quad \forall A \in \mathcal{E}, \quad Y_n^{-1}(A) \in \mathcal{F}_{n-1} $$
        par convention, on dira que $ \mathcal{F}_{-1} = \{ \emptyset, \Omega\}$. 
    \end{enumerate}
    La notion de \emph{prévision} signifie que l'on peut connaître $Y_n$ avant le temps $n$ grâce 
    aux informations contenues dans $ \mathcal{F}_{n_1}$. 
\end{definition}

La notion de filtration permet aussi de généraliser les temps d'arrêts. 

\begin{definition}[Temps d'arrêt]
    Un temps d'arrêt $T$ pour une filtration $ \mathcal{F}_n$ sur $( \Omega, \mathcal{A}, \myP)$ est une 
    variable aléatoire à valeurs dans $( \overline{\N}, \mathcal{P}( \overline{\N}))$ telle que : 
        \begin{align*}
            \forall n \in \ N, \quad  & \{T \leqslant n\} \in \mathcal{F}_n \\  
        \end{align*} 
    Pour être plus précis, cela revient à $ T^{-1} \{0, 1, 2, \dots, n\} = \{ \omega \in \Omega \mid T( \omega) \leqslant n\} \in \mathcal{F}_n$. 
\end{definition}
